\section{Наивная теория вероятностей.}
\paragraph{Схема Бернулли.}
Будем называть схемой Бернулли эксперимент бросания несимметричной монетки счётного числа раз: вероятность орла -- $p$, решки $q = 1 - p$ в каждом бросании.
Надо найти распределение случайной величины $\xi$, равной номеру первого успеха.
Очевидно, что $\P(\xi = N) = q^{N - 1}p$.
Такое распределение называется геометрическим.
\paragraph{Биномиальное распределение.}
Теперь найдём распределение вероятность того, что из $N$ испытаний будет ровно $n$ успехов:
\[
    \P(\eta = n) = \binom{N}{n}p^{n}q^{N - n}.
\]
Такая $\eta \sim Binom(N, p)$.
\paragraph{Предельный переход к распределению Пуассона.}
Предположим теперь, что в серии из $N$ бросаний ровно $n$ успехов и при этом $p \ra 0, N \ra +\infty$ и $Np = \lambda$.
Производящей функцией биномиального распределения является
\[
    G(z) = (q + pz)^N.
\]
При таком предельном переходе мы получаем
\[
    G(z) = (q + \dfrac{\lambda}{N}z)^N = (1 + \dfrac{\lambda(z - 1)}{N})^N \ra e^{\lambda (z - 1)} = \sum\limits_{k = 0}^{+\infty} \dfrac{(\lambda z)^k}{k!}e^{-\lambda}, N \ra \infty.
\]
